\section{Machines, clocks, and resource accounting}
\label{sec:machines}

\subsection{A concrete bounded computation predicate}

For time we use ordinary deterministic single-tape machines, with one fixed
global alphabet \(\{0,1,\#,\mathrm{blank}\}\), a finite state set, an initial state,
and distinct accepting and rejecting states. The tape is indexed by integers;
the input occupies positions \(0,\ldots,n-1\), the remaining cells are blank,
and the head starts at 0. A transition reads a symbol, changes state, writes
a symbol, and moves by \(-1,0\), or 1. Halting states are made absorbing,
with no change to the tape or head. Undefined or malformed transitions reject.
A malformed machine code denotes a fixed rejecting machine. Numerical code
0 for a machine description is normalized to the empty binary description
(word code 1), which is interpreted as rejecting. Thus all coded
operations are total numerical functions.

A configuration consists of a state, an integer head position, and a finite
record of the nonblank tape cells. Encode a signed position by its sign and
absolute value, and normalize the tape record by removing blank boundary
cells. Reading an unrecorded cell returns blank. Lookup, replacement, adding
one new boundary cell, and normalization are bounded finite loops on this
record. Consequently \(\Step(e,C)\) and the initial configuration
\(\operatorname{Init}(e,x)\) are primitive recursive uniformly in the machine
description \(e\). Define
\[
 \Run(e,x,0)=\operatorname{Init}(e,x),\qquad
 \Run(e,x,t+1)=\Step(e,\Run(e,x,t)),
\]
and let \(\Acc(e,x,t)\) be 1 precisely when the resulting state is accepting.
For invalid input code 0 it is defined to be 0.

The equality \(\Acc(e,x,t)=1\) means ``accepts by time \(t\).'' It does not
assert that the original unclocked machine always halts. A clocked evaluator
runs the indicated computation and returns 0 on timeout, thereby always
producing a bit.

\begin{lemma}[Bounded computation]
\label{lem:boundedrun}
The functions \(\Run\) and \(\Acc\) are primitive recursive. Their defining
equations, trace composition, uniqueness of a checked deterministic trace,
and monotonicity of accepting-by-time are provable in \ipra.
\end{lemma}

\begin{proof}
Iteration is primitive recursion on \(t\), with \(e,x\) as parameters.
After \(t\) steps the head lies in \([-t,t]\), and all visited or initially
nonblank cells lie in \([-t,n+t]\). A window of
\(W=n+2t+3\) cells therefore contains the entire relevant tape, including
blank sentinels. For a fixed finite alphabet of size \(g\) and \(q\) states,
a window configuration is bounded, for example, by
\(qWg^W\), after choosing an offset for the head. Uniformly in machine codes,
the decoded numbers \(q,g\) provide the same primitive-recursive expression.

For composition prove
\(\operatorname{Iter}_{e}(r+s,C)=
\operatorname{Iter}_{e}(s,\operatorname{Iter}_{e}(r,C))\)
by induction on \(s\), using equality of coded configurations as the
induction formula. A trace checker tests its initial row and each adjacent
pair by a bounded conjunction. Inducting on the row number proves that every
checked row equals the corresponding iterate; the prefix assertion is a
finite conjunction and is Boolean-compiled. Finally an accepting
configuration is a fixed point of \(\Step\), so iteration preserves it.
This gives accepting-by-time monotonicity.
\end{proof}

For \(T\) steps, an entire window trace can be bounded by
\(B^{T+1}\), with \(B=1+qWg^W\), and can be generated by a single loop
accumulating rows. Such a trace is useful for correctness certificates; it is
not necessarily the storage used by the machine being analyzed.
When \(\Run\) uses the normalized finite-support representation, first
convert each row to the fixed window selected by \(T\). Bounded cell lookup
defines this conversion; the head/window bound proves that applying a step
then converting agrees with the fixed-window step throughout the trace.
The bound \(B^{T+1}\) applies to these converted rows, not to an unrelated
sparse-record serialization.

\subsection{Physical time is a second proof}

The arithmetic definition of \(\Run\) proves totality of the clocked
operation. It does not yet prove a feasible simulation. We use the following
coarse implementation bound; the later diagonal section gives its sweep
construction explicitly.

\begin{lemma}[Coarse universal implementation]
\label{lem:coarseU}
There is a fixed ordinary machine and a fixed positive integer \(C\) which
compute \(\Acc(e,x,T)\), with a unary clock supplied, in at most
\[
 C(\len e+\len x+T+1)^4
\]
steps. The transition-table simulation, its correctness, and the polynomial
overhead bounds have bounded invariants and \ipra{} proofs.
\end{lemma}

The constant \(C\) depends on the chosen universal implementation; this
package gives a constructive implementation recipe, not a numerically
optimized transition table for that universal machine. The full Lean
formalization must realize this recipe and establish its cost lemma. The
degree four is a deliberately generous bound for fixed-track sweep
simulation and table lookup, not a sharp time-hierarchy claim.

Fixed finitely many tapes can be represented on tracks of one tape, with
marked head positions. A complete sweep reads the marked symbols, and a
second updates the tracks and their markers. The resulting larger, but
fixed, tracking alphabet is encoded into the chosen global alphabet by
constant-length blocks, adding only a fixed multiplicative overhead.
At time \(t\), the occupied
interval has length at most a fixed multiple of \(n+t+1\). Thus \(T\)
multitape moves can be simulated in polynomial time. Conversely, one tape is
a special case of finitely many tapes. This supplies the model invariance
needed for polynomial and exponential classes. We invoke neither a
Church--Turing thesis nor a unit-cost model of arbitrary-precision arithmetic
as a resource theorem.

\subsection{Clock families and asymptotic bounds}

Write
\[
 p_{c,k}(n)=c(n+1)^k\quad(c\ge1,\ k\ge0),\qquad
 e_{c,k}(n)=2^{c(n+1)^k}.
\]
The evaluations are primitive recursive jointly in \(c,k,n\). Coefficients
and exponents are numerical data, with no arbitrary-function evaluation.
Whenever input reading is needed, enlarge a bound to dominate \(n+1\).
Every polynomial with nonnegative integer coefficients is bounded by some
\(p_{c,k}\). Sums and compositions have explicit majorants: for
\(p=a(n+1)^r\), \(q=b(n+1)^s\), with \(r,s\ge1\),
\[
 p(n)+q(p(n))+n+1
 \le\bigl(a+b(a+1)^s+1\bigr)(n+1)^{\max(r,rs,1)}.
\]
This follows by \(p(n)+1\le(a+1)(n+1)^r\), multiplication, and monotonicity
of powers. Any fixed polynomial simulation overhead can then be absorbed by
multiplying the degree and enlarging the coefficient.

It is useful to encode a polynomial as a finite list of nonnegative
coefficients as well. Its evaluation, addition, multiplication, and
substitution are primitive recursive, and their identities are proved by
finite sums and products. The Lean core implements an equivalent syntax of
constants, the variable, addition, and multiplication. The choice between
these polynomial representations is an implementation choice; it introduces
no new arithmetic strength.

Big-O notation can be expressed by explicit constants and a threshold.
For a known clocked algorithm, finitely many lengths below the threshold can
be covered by a bounded maximum and a larger coefficient. We normally
present bounds valid at every length, which avoids an unnecessary
eventual-domination obligation. For an arbitrary potentially partial machine,
one cannot form its worst-case halting time without a supplied halting bound.

\subsection{Separate conventions for space and output}

Space classes use a read-only input tape with endmarkers and a fixed number
of one-way work tapes. The input cells are not charged as work space; input
head positions still contribute to configuration counts. A cap \(s\ge1\)
allows cells \(0,\ldots,s-1\) on each work tape. An attempted transition
outside the cap has no accepting continuation. Section~\ref{sec:space}
constructs the finite configuration graph and proves cutoff and simulation
theorems. These explicit conventions are necessary for logarithmic space.

A function machine has a separate initially blank output tape, appending at
most one output bit per charged transition. Its output length is therefore
at most its clock, up to fixed record headers. A clocked transducer returns
the empty word on timeout; when used as a reduction, its entire clocked
behavior is the specified function. Output-size and runtime bounds are
stored together so that a later algorithm has a justified input-size bound.
